\section{Related work}\label{sec:related}

\paragraph{Tractability of sparse nonconvex quadratic programs.}
For binary quadratic and pseudo-Boolean optimization, bounded treewidth gives
polynomial algorithms by nonserial dynamic programming
\cite{BerteleBrioschi1972,CramaHansenJaumard1990,Dechter1999} and exact
lift-and-project or linear formulations of size exponential only in the width
\cite{BienstockOzbay2004,WainwrightJordan2004,Laurent2009}; sparse moment and
sum-of-squares hierarchies exploit the same structure for polynomial problems
\cite{WakiEtAl2006,Lasserre2006}. For continuous box QP, the tractable
classes of unbounded treewidth of Del Pia and Khajavirad
\cite{DelPiaKhajavirad2026} and of Khajavirad \cite{Khajavirad2026PolyBox}
take variables with nonpositive diagonal to be binary
(Remark~\ref{rem:cv-instances} compares Khajavirad's elimination of the
continuous components with our value factors). The exact
piecewise-quadratic messages of Del Pia and Khajavirad's forest algorithm,
also used for
total-variation problems on trees
\cite{KolmogorovPockRolinek2016,KuricAhmetspahicPock2024}, are compared with
our grids in Section~\ref{sec:limits-messages}. Box QP is also polynomial for
fixed Hessian rank \cite{HladikCernyRada2021}. For convex quadratics with
indicator variables, Bhathena et al.\ give dynamic programs on trees
\cite{BhathenaEtAl2026Trees} and on graphs of bounded treewidth, the latter
linear in $n$ for fixed width, margin, volume growth and condition numbers
\cite{BhathenaEtAl2026Graphs}; this is the closest treewidth-and-conditioning
result we know (see also \cite{BienstockChen2024}). For integer QP,
bounded treewidth together with bounded domains gives fixed-parameter
tractability \cite{EibenEtAl2019}; parameterized by the number of variables,
the problem is fixed-parameter tractable with bounded coefficients
\cite{Lokshtanov2015} and W[1]-hard without them \cite{Herrmann2026}; see \cite{GanianOrdyniak2018} for integer linear
programming. In fixed dimension, integer QP in the plane
\cite{DelPiaWeismantel2014} and, according to a recent preprint,
mixed-integer QP (MIQP)
\cite{AriHildebrand2026} are solved exactly in polynomial time. Approximation
schemes polynomial in $1/\varepsilon$ exist in fixed dimension
\cite{DeLoeraEtAl2008,HildebrandWeismantelZemmer2016}, for sparse polynomial
problems with constraints \cite[Theorem~4]{BienstockMunoz2018}, for QP with a
fixed number of negative eigenvalues in exact real arithmetic
\cite{Vavasis1992}, and for MIQP with fixed rank and a fixed number of
integer variables \cite{DelPia2023} or, in the bit model, a fixed number of
negative eigenvalues and of integer variables \cite{DelPia2026Jacobi}. Several of these works explain why, in their
settings, the dependence on $1/\varepsilon$ cannot be polylogarithmic unless
P${}={}$NP \cite{BienstockMunoz2018,Vavasis1992,DelPia2026Jacobi}, from the
NP-hardness of Subset Sum or of QP with one negative eigenvalue
\cite{PardalosVavasis1991}. Running times that depend on a condition number
of the instance are standard in numerical computation
\cite{BurgisserCucker2013}.

\paragraph{Grid and vertex bounds, filtering and certificates.}
Lower bounds from function values on a grid are classical for Lipschitz
functions \cite{Piyavskii1972,Shubert1972,Nesterov2004}; with second-order
information the error is quadratic in the mesh
\cite{BreimanCutler1993,DeKlerkLasserreLaurentSun2017}. Separable quadratic
perturbations that convexify a function ($\alpha$BB) have maximum separation
$\frac14\sum_i\alpha_i(u_i-\ell_i)^2$
\cite{MaranasFloudas1994,AndroulakisMaranasFloudas1995,AdjimanEtAl1998};
perturbations of the opposite sign make it edge-concave, with a
vertex-polyhedral envelope
\cite{Tardella2004,MeyerFloudas2005,Hasan2018,MisenerFloudas2012}, and Bajaj
and Hasan evaluate such an underestimator, built from an upper bound on the
diagonal Hessian entries, at the $2^n$ vertices of a box
\cite{BajajHasan2020}. Piecewise-linear and sawtooth relaxations of squares
apply the same second-order error termwise inside a mixed-integer program
(MIP) \cite{BeachHildebrandHuchette2022,BeachEtAl2024,DongLuo2018}.

Removing a region whose relaxation value exceeds an incumbent is
optimality-based bound tightening and branch-and-reduce
\cite{RyooSahinidis1996,BelottiEtAl2009,GleixnerEtAl2017,PuranikSahinidis2017},
cost-based domain filtering in constraint programming
\cite{FocacciLodiMilano1999} and dead-end elimination in protein design
\cite{DesmetEtAl1992,GainzaRobertsDonald2012}; thresholded min-marginals,
computed by two passes of message passing
\cite{Dawid1992,AjiMcEliece2000,WainwrightJaakkolaWillsky2005,KohliTorr2006},
prune states in structured prediction cascades \cite{WeissTaskar2010}. Exact
rational MIP \cite{CookKochSteffyWolter2013,EiflerGleixner2023} and
independently checkable branch-and-bound certificates
\cite{CheungGleixnerSteffy2017} guard MIP results against floating-point
error, and our path certificate follows the latter, with a tree that is a
path; interval methods and safe bounds control rounding errors instead
\cite{Hansen1980,HansenWalster2004,NeumaierShcherbina2004}. Rational
sum-of-squares certificates give exact lower bounds for polynomial objectives
\cite{PeyrlParrilo2008,KaltofenEtAl2008}, but these bounds are values of
relaxations, whereas a path certificate brackets $\OPT$ within any prescribed
$2^{-q}$.

\paragraph{Accuracy-independent work under growth.}
With second-order convergent bounds, the number of branch-and-bound boxes
near a nondegenerate minimizer is independent of the tolerance (the cluster
problem), but local, asymptotic and possibly exponential in the dimension
and in the conditioning of the Hessian
\cite{DuKearfott1994,Neumaier2004,WechsungSchaberBarton2014,KannanBarton2017}.
Lemmas~\ref{lem:inv} and~\ref{lem:states} are a global, nonasymptotic form of
this phenomenon in which the exponential dependence falls on the bag size,
and $8\bar\kappa\theta^2\le1$ plays the role of the bound on the second-order
prefactor. Optimistic optimization attains exponential rates when the growth
order at the optimum matches the assumed smoothness \cite{Munos2011}; for
Lipschitz functions with small near-optimal sets, branch-and-bound needs
$O(\log(1/\varepsilon))$ function values \cite{Perevozchikov1990}, and
Bachoc, Cesari and Gerchinovitz characterize, instance by instance, the
number of evaluations needed to find and certify an $\varepsilon$-optimal
point \cite{BachocCesariGerchinovitz2021}. Quadratic growth and error bounds
give linear convergence of local methods
\cite{LuoTseng1993,KarimiNutiniSchmidt2016,DrusvyatskiyLewis2018,
NecoaraNesterovGlineur2019}, and restarts remove the need to know the growth
constant \cite{RouletDAspremont2020}, as our capped trials do.

Proximity and scaling give algorithms logarithmic in the domain size for
linear and separable convex integer programs
\cite{CookGerardsSchrijverTardos1986,HochbaumShanthikumar1990,
GranotSkorinKapov1990,HunkenschroderEtAl2026,EisenbrandEtAl2025}. The closest
classical analogue of our stages is the scaling algorithm of Hochbaum and
Shanthikumar. For a separable convex objective over $\{x:Ax\ge b\}$, each
iteration minimizes a piecewise-linear interpolation of the objective on a
grid of $O(n\Delta_A)$ points per variable, where $\Delta_A$ bounds the
absolute values of the subdeterminants of $A$. The grid lies in a box around
the previous solution, and a proximity theorem allows the box to be halved.
For continuous variables, $\lceil\log_2(s/(2\varepsilon))\rceil$ such
grids, with $s$ the initial box width, give a point within $\varepsilon$ of
an optimal solution in the maximum norm, guaranteed but not certified
\cite[Section~1.3 and Theorems~1.1, 3.8 and~4.4]{HochbaumShanthikumar1990}.
For totally unimodular $A$ and integral data, the integer problem is solved
by one linear program over the interpolation at all integer points
\cite{Meyer1977}, or by one such linear program per scale
\cite[Theorem~4.3]{HochbaumShanthikumar1990}; Lemma~\ref{lem:tu-round}
uses the same integrality. Our objectives are nonseparable and nonconvex. The
retained region comes from min-marginal filtering, which needs no assumption
and yields a certificate; growth only bounds its size. The grid problems are
solved by dynamic programming over a tree decomposition.

\paragraph{Adaptive discretization.}
Coarse-to-fine dynamic programming refines superstates with optimistic costs
along the current optimal path, so that each coarse problem is a lower bound
\cite{Raphael2001,FelzenszwalbMcAllester2007}. Interval methods give discrete
lower bounds for continuous Markov random fields \cite{PengEtAl2011} and, for
optimal power flow on trees, a number of bound evaluations polynomial in
$1/\varepsilon$ \cite[Theorem~2]{DvijothamEtAl2017}. Decision diagrams give
outer approximations of integer nonlinear programs \cite{DavarniaVanHoeve2021}
and, on partitions of continuous domains with spatial branch-and-bound, a
global method for bounded mixed-integer nonlinear programs
\cite{DavarniaKiaghadiQiu2026}. Continuous distributed constraint optimization
\cite{HoangEtAl2020,TroullinosEtAl2022}, variable-resolution control
\cite{MunosMoore2002}, corridor methods \cite{HeidariEtAl1971}, two-level
dynamic programming \cite{LeaverFayKuhlmanSnoeyink2005}, adaptive
multivariate partitioning with bound tightening \cite{NagarajanEtAl2019} and
re-centered discretizations of fixed size \cite{GupteKosterKuhnke2022} also
refine discretizations adaptively. None of these methods certifies a given
accuracy with a number of states independent of that accuracy.

\paragraph{Exact rational output.}
A rational QP that attains its minimum has a global minimizer of polynomial
encoding length \cite{Vavasis1990,DelPiaDeyMolinaro2017}; the
stationary-polytope argument of Section~\ref{sec:exact} follows Vavasis, and
Remark~\ref{rem:np} discusses NP membership. Rounding to an exact rational is classical
\cite[Section~5.1]{GrotschelLovaszSchrijver1988},
\cite{KozlovTarasovKhachiyan1980}, and in linear programming the optimal face
is identified once the gap is below a threshold set by the encoding length
\cite{Ye1992,MehrotraYe1993}; Lemma~\ref{lem:snap} does this for box QP, with
set growth turning a small gap into proximity to the optimal set. Uniqueness
of a QP minimizer implies quadratic growth \cite{Contesse1980}, and growth
towards the optimal set of a QP over a polytope follows from the error bound
of Luo and Sturm \cite[Theorem~3.3]{LuoSturm2000}. For polynomials of degree
three or more, irrational minimizers and encoding-size phenomena limit exact
output \cite{BienstockDelPiaHildebrand2023}.

\paragraph{Value functions and multiparametric programming.}
Eliminating private variables through their optimal response is standard in
Benders decomposition \cite{Benders1962,Geoffrion1972} and in condensing for
model predictive control
\cite{BockPlitt1984,JerezKerriganConstantinides2012,Axehill2015}.
Multiparametric quadratic programming gives piecewise-affine minimizers and
piecewise-quadratic values on polyhedral critical regions
\cite{BemporadEtAl2002,TondelJohansenBemporad2003}, also with nonunique
minimizers \cite{SpjotvoldTondelJohansen2007,PatrinosSarimveis2011}; value
functions of parameters that enter the objective linearly over a fixed
feasible set are
concave \cite{MangasarianRosen1964,FiaccoKyparisis1986}, and infima preserve
semiconcavity \cite{CannarsaSinestrari2004}. We use these facts for a global
upper coordinate curvature bound of the reduced objective. Low negative
inertia with convex recourse is treated by spatial branching in
\cite{Vavasis1992,LuoBaiLimPeng2019}.

\paragraph{Cuts and submodularity.}
Binary quadratic functions with nonpositive pair coefficients are minimized by
one minimum cut \cite{Ivanescu1965,PicardRatliff1975,BorosHammer2002,
KolmogorovZabih2004}, also after sign switching when the signed interaction
graph is balanced \cite{Harary1953,HammerHansenSimeone1984,Padberg1989}.
Variables with nonpositive diagonal can be fixed to endpoints
\cite{Rosenberg1972,DelPiaKhajavirad2026,Khajavirad2026PolyBox}, so continuous
quadratics that are submodular and concave in each coordinate reduce to such
cuts \cite{BurerNatarajanWillemsen2026v3}, and pairwise problems submodular
with respect to label orders reduce to cuts on threshold encodings
\cite{SchlesingerFlach2006,Ishikawa2003,Hochbaum2001}. For general continuous
submodular quadratics, Burer, Natarajan and Willemsen show that a
semidefinite relaxation is exact up to dimension three and give a
four-variable gap; whether the problem is polynomial when the dimension is
part of the input is open \cite[footnote~2]{BurerNatarajanWillemsen2026v3}
(for each fixed dimension it is, by enumerating faces \cite{Vavasis1990}). A
sign-restricted case is exact \cite{KimKojima2003}, and discretization
methods \cite{Bach2019,Bach2018Isotonic,AxelrodLiuSidford2020} are polynomial
in $1/\varepsilon$; for balanced box quadratics under point growth,
Theorem~\ref{thm:balanced} is polynomial in $I$, $\kappa$ and
$\log(1/\varepsilon)$, which does not settle
the open question because $\kappa$ can be exponential in $I$. Partial minimization preserves
submodularity \cite{Topkis1978,Topkis1998}, which underlies the mixed
continuous--discrete methods of \cite{GomezHan2025,BuntonTabuada2022}, with
certificates from greedy extreme bases
\cite{Edmonds1970,Cunningham1985,IwataFleischerFujishige2001}.

\paragraph{Constraints, optimal sets and limits.}
Integrality of totally unimodular fibers is the Hoffman--Kruskal theorem
\cite{HoffmanKruskal1956,Schrijver1986}, and marginal-preserving rounding
under hard constraints is classical
\cite{GandhiEtAl2006,ChekuriVondrakZenklusen2010}; Remark~\ref{rem:tu-bm}
compares Section~\ref{sec:constraints} with the linear programs of Bienstock
and Mu\~noz \cite{BienstockMunoz2018}. Section~\ref{sec:endpointset} relates
our description of the optimal set to Rosenberg's face criterion and to tree
reparameterizations
\cite{Rosenberg1972,WainwrightJaakkolaWillsky2005,Werner2007}. Diagonal
Lagrangian certificates for box QP are given in
\cite{BeckTeboulle2000,JeyakumarRubinovWu2006,LiWuQuan2015}, and exactness
of relaxations of box QP is characterized in \cite{QiuYildirim2024}; sparse
second-order cone and semidefinite formulations are exact under graph
conditions \cite{SojoudiLavaei2014,DeyKhajavirad2026,Khajavirad2026SparseSDP},
and chordal sparsity guarantees positive semidefinite completions
\cite{GroneEtAl1984,VandenbergheAndersen2015}. Marginals that agree on
separators glue into a joint measure on tree-structured families
\cite{Vorobev1962,WainwrightJordan2008}, but locally consistent moments need
not \cite{SheraliAdams1990,WainwrightJordan2008}, and complete parametric
descriptions can have exponentially many pieces
\cite{Murty1980,MairalYu2012,GartnerJaggiMaria2012}.
